\documentclass[12pt]{article}
\usepackage[utf8]{inputenc}
\usepackage[T1]{fontenc}
\usepackage{hyperref}
\title{Northwestern !Kung Language Resources for K-12 Curriculum Development}
\author{Compiled by Claude Sonnet 5 (claude-sonnet-5) for Andr\'as Kornai}
\date{}
\begin{document}
\maketitle
\section{Introduction}
Northwestern !Kung (ISO 639-3 vaj, Glottocode vase1234) is a Kx'a language of Angola and Namibia. Caution: the search results mostly concern related varieties, namely Ekoka !Kung (ISO code knw, per its Wikipedia article) and Ju$|$'hoan, which are separate codes. I found nothing specific to vaj, and I could not confirm that the Ekoka material applies to it.
\section{Primary Educational Resources}
None found for vaj. Wikipedia states that Ju$|$'hoan is the only !Kung variety with a written tradition, and that it is not sufficiently intelligible with the Northwestern dialects for the same literature to be used.
\section{Linguistic and Reference Materials}
\begin{itemize}
\item Ekoka !Kung, Wikipedia (related variety): \url{https://en.wikipedia.org/wiki/Ekoka_%C7%83Kung}
\item !Kung languages, Wikipedia: \url{https://en.wikipedia.org/wiki/%C7%83Kung_languages}
\item Tsammalex language page (wordlist after K\"onig \& Heine 2008): \url{https://tsammalex.clld.org/languages/knw}
\end{itemize}
\section{Religious and Community Resources}
None found.
\section{Digital and Audio Resources}
None found for vaj.
\section{Curriculum Development Context}
No orthography or school materials were found. A curriculum would start from documentation work; K\"onig \& Heine (2008) and Heine \& Honken (2010) were named as leads, but I did not see them directly.
\end{document}